% ============================================================
% 第6章 原子结构与元素性质（选必2第一章）
% ============================================================
\chapter{原子结构与元素性质}

\begin{daoxue}
\textbf{知识主线：}电子在原子核外怎样排布（能层、能级、原子轨道、构造原理、三条排布规则）$\to$ 排布决定元素在周期表中的位置（分区）$\to$ 位置决定性质的周期性变化（原子半径、电离能、电负性）$\to$ 用“结构—位置—性质”推断元素。

\textbf{高考考情：}广东卷选择题每年有一道元素推断（2024年第12题由结构式推断、2025年第9题由第一电离能图推断、2026年第11题由轨道“半充满”推断）；大题中常考价层电子排布式或轨道表示式（2025年第19题(1) Ti、2026年第19题(1) S）。

\textbf{学习目标：}①正确书写 1～36 号元素的电子排布式、价层电子排布式和轨道表示式；②理解电离能、电负性的递变规律及特例；③熟练运用“位—构—性”推断元素。
\end{daoxue}

\section{原子核外电子排布}

\subsection{能层、能级与原子轨道}
核外电子按能量不同分成\textbf{能层}（$n=1,2,3,\cdots$，符号 K、L、M、N…），每个能层又分为若干\textbf{能级}（s、p、d、f），能层序数 $n$ 就是该能层的能级数。电子在核外空间出现的概率分布用\textbf{电子云}描述，人们把电子出现概率约 90\% 的空间圈出来，称为\textbf{原子轨道}：s 轨道呈球形（每个 s 能级 1 个轨道），p 轨道呈哑铃形（每个 p 能级 3 个相互垂直的轨道 $p_x$、$p_y$、$p_z$），d 能级有 5 个轨道，f 能级有 7 个轨道。每个原子轨道最多容纳 2 个自旋相反的电子，所以 s、p、d、f 能级最多容纳 2、6、10、14 个电子，第 $n$ 能层最多容纳 $2n^2$ 个电子。

\subsection{构造原理与排布规则}
\begin{gainian}[构造原理与三条规则]
\textbf{构造原理}（能级填充顺序）：1s 2s 2p 3s 3p \textbf{4s 3d} 4p 5s 4d 5p 6s 4f 5d 6p…（出现“能级交错”：4s 能量低于 3d，先填 4s）。\\
\textbf{能量最低原理}：基态原子电子优先占据能量低的轨道，使整个原子能量最低。\\
\textbf{泡利原理}：一个原子轨道最多容纳 2 个电子，且自旋相反。\\
\textbf{洪特规则}：电子填入同一能级的不同轨道时，优先\textbf{单独}分占不同轨道，且自旋平行。特例：能级处于全满（$p^6$、$d^{10}$）、半满（$p^3$、$d^5$）或全空时较稳定，如 Cr：[Ar]$3d^54s^1$，Cu：[Ar]$3d^{10}4s^1$。
\end{gainian}

\begin{center}
\small
\begin{tabular}{lll}
\toprule
表示方法 & 以 Fe（26 号）为例 & 说明 \\
\midrule
电子排布式 & $1s^22s^22p^63s^23p^63d^64s^2$ & 书写时按能层顺序，3d 写在 4s 前 \\
简化排布式 & [Ar]$3d^64s^2$ & 用稀有气体结构代替内层 \\
价层电子排布式 & $3d^64s^2$ & 主族元素为最外层，过渡元素含次外层 d \\
价层轨道表示式 & 3d：\fbox{$\uparrow\downarrow$}\fbox{$\uparrow$}\fbox{$\uparrow$}\fbox{$\uparrow$}\fbox{$\uparrow$}\ 4s：\fbox{$\uparrow\downarrow$} & 4 个未成对电子 \\
离子 & \ce{Fe^2+}：[Ar]$3d^6$；\ce{Fe^3+}：[Ar]$3d^5$ & \textbf{先失最外层（4s）电子} \\
\bottomrule
\end{tabular}
\end{center}

\begin{benzhi}[为什么 \ce{Fe^3+} 比 \ce{Fe^2+} 稳定，\ce{Cu+} 在固态下比 \ce{Cu^2+} 稳定]
\ce{Fe^3+} 的 3d 能级为 $3d^5$ 半充满，较稳定，所以 \ce{Fe^2+}（$3d^6$）易被氧化为 \ce{Fe^3+}。\ce{Cu+} 为 $3d^{10}$ 全充满，所以高温下 CuO 分解生成 \ce{Cu2O}。但在水溶液中 \ce{Cu+} 易歧化（\ce{2Cu+ = Cu + Cu^2+}），因为 \ce{Cu^2+} 的水合放出更多能量——判断稳定性还要看具体环境。
\end{benzhi}

\begin{yicuo}
\begin{enumerate}
\item 轨道表示式中违反泡利原理（同一方框内两电子自旋相同）或违反洪特规则（未分占、或自旋不平行）都是错误的。
\item 书写过渡元素阳离子的电子排布时，不能先失去 3d 电子。
\item 基态原子“未成对电子数”要画轨道表示式数：第四周期 Cr 有 6 个（$3d^54s^1$），是第四周期最多的。
\end{enumerate}
\end{yicuo}

\begin{liti}[电子排布的书写]
(1) 写出基态 Ga、Br 原子的简化电子排布式；(2) 写出 \ce{Cu^2+}、\ce{Mn^2+} 的价层电子排布式；(3) 基态 Co、Ni 原子的未成对电子数分别是多少？(4) 第四周期中，基态原子未成对电子数与 Fe 相同的元素有哪些？
\tcblower
\jie (1) Ga：[Ar]$3d^{10}4s^24p^1$（不要漏掉 $3d^{10}$）；Br：[Ar]$3d^{10}4s^24p^5$。(2) \ce{Cu^2+}：$3d^9$（Cu 为 $3d^{10}4s^1$，先失 4s 上 1 个电子，再失 3d 上 1 个）；\ce{Mn^2+}：$3d^5$。(3) Co：$3d^74s^2$，3d 上 7 个电子，5 个轨道中 2 个成对、3 个单电子，未成对电子数 3；Ni：$3d^84s^2$，未成对电子数 2。(4) 没有。Fe（$3d^64s^2$）有 4 个未成对电子；未成对电子数为 4 只可能是 $3d^4$ 或 $3d^6$ 构型，而 24 号 Cr 因半满稳定为 $3d^54s^1$（6 个未成对电子），所以第四周期中只有 Fe 有 4 个未成对电子。
\end{liti}

\subsection{原子光谱}
基态原子吸收能量，电子跃迁到较高能级成为激发态；激发态电子跃迁回较低能级时以光的形式释放能量。不同元素的原子发生跃迁时吸收或发射的光的频率不同，形成特征的\textbf{原子光谱}，可用于鉴定元素（光谱分析）。\textbf{焰色}就是电子由激发态跃迁到低能级时放出的光，属于发射光谱，是物理变化。

\section{原子结构与元素周期表}
\textbf{周期}：周期序数$=$能层数；每一周期从 $ns^1$ 开始到 $ns^2np^6$ 结束（第一周期到 $1s^2$）。\textbf{族}：主族序数$=$价层电子数；第ⅢB～ⅦB族序数$=$价层电子（$(n-1)d+ns$）数。\textbf{分区}：

\begin{center}
\small
\begin{tabular}{llll}
\toprule
分区 & 范围 & 价层电子排布 & 特点 \\
\midrule
s 区 & ⅠA、ⅡA（含 H、He） & $ns^{1\sim2}$ & 活泼金属（H 除外） \\
p 区 & ⅢA～ⅦA、0 族 & $ns^2np^{1\sim6}$（He 除外） & 包括全部非金属（H 除外） \\
d 区 & ⅢB～ⅦB、Ⅷ族 & $(n-1)d^{1\sim9}ns^{1\sim2}$（Pd 例外） & 过渡金属，多种化合价 \\
ds 区 & ⅠB、ⅡB & $(n-1)d^{10}ns^{1\sim2}$ & Cu、Zn 等 \\
f 区 & 镧系、锕系 & $(n-2)f\,(n-1)d\,ns$ & 性质相似 \\
\bottomrule
\end{tabular}
\end{center}

\section{元素周期律}

\subsection{原子半径}
\textbf{同周期}从左到右，核电荷数增大而能层数不变，核对外层电子的吸引增强，原子半径减小；\textbf{同主族}从上到下，能层数增加，原子半径增大。\textbf{离子半径}：电子层结构相同的离子，核电荷数越大，半径越小，如 $r(\ce{O^2-})>r(\ce{F-})>r(\ce{Na+})>r(\ce{Mg^2+})>r(\ce{Al^3+})$；同种元素 $r(\text{阴离子})>r(\text{原子})>r(\text{阳离子})$，$r(\ce{Fe^2+})>r(\ce{Fe^3+})$。

\subsection{电离能}
\begin{gainian}[第一电离能 $I_1$]
气态电中性基态原子失去一个电子转化为气态基态正离子所需的最低能量。$I_1$ 越小，越容易失电子，金属性越强。同一原子逐级电离能 $I_1<I_2<I_3<\cdots$。
\end{gainian}

\textbf{递变规律：}同周期从左到右 $I_1$ 呈增大趋势，同主族从上到下减小。\textbf{特例}：ⅡA 族（$ns^2$ 全满）和 ⅤA 族（$np^3$ 半满）的 $I_1$ 大于同周期相邻元素。如第二周期：Li$<$B$<$Be$<$C$<$O$<$N$<$F$<$Ne；第三周期：Na$<$Al$<$Mg$<$Si$<$S$<$P$<$Cl$<$Ar。

\begin{center}
\begin{tikzpicture}
\begin{axis}[width=10cm,height=4.6cm,xlabel={原子序数},ylabel={$I_1$/(kJ·mol$^{-1}$)},xmin=0,xmax=21,ymin=0,ymax=2600,
 xtick={1,2,...,20},tick label style={font=\scriptsize},label style={font=\small},grid=major,grid style={gray!20}]
\addplot[thick,mainblue,mark=*,mark size=1.4pt] coordinates {(1,1312)(2,2372)(3,520)(4,900)(5,801)(6,1086)(7,1402)(8,1314)(9,1681)(10,2081)(11,496)(12,738)(13,578)(14,787)(15,1012)(16,1000)(17,1251)(18,1521)(19,419)(20,590)};
\node[font=\scriptsize,above] at (axis cs:2,2372) {He}; \node[font=\scriptsize,above] at (axis cs:10,2081) {Ne}; \node[font=\scriptsize,above] at (axis cs:18,1521) {Ar};
\node[font=\scriptsize,below] at (axis cs:3,520) {Li}; \node[font=\scriptsize,below] at (axis cs:11,496) {Na}; \node[font=\scriptsize,below] at (axis cs:19,419) {K};
\node[font=\scriptsize,above] at (axis cs:4,900) {Be}; \node[font=\scriptsize,above] at (axis cs:7,1402) {N}; \node[font=\scriptsize,above] at (axis cs:12,738) {Mg}; \node[font=\scriptsize,above] at (axis cs:15,1012) {P};
\end{axis}
\end{tikzpicture}
\end{center}
\noindent\small 图：1～20 号元素的第一电离能。每一周期从碱金属（谷）到稀有气体（峰）总体增大；ⅡA（Be、Mg）、ⅤA（N、P）出现“凸起”。2025年广东卷第9题正是据此图形推断元素。\normalsize

\textbf{逐级电离能的应用}：若 $I_{n+1}\gg I_n$（出现突跃），说明失去 $n$ 个电子后达到稳定结构，该元素最外层有 $n$ 个电子，常见化合价为 $+n$。如某元素 $I_1\sim I_4$（kJ·mol$^{-1}$）为 578、1817、2745、11577，$I_4$ 突跃，为 ⅢA 族的 Al。

\begin{liti}[由电离能推断元素]
短周期元素 X、Y 的逐级电离能（kJ·mol$^{-1}$）如下：X：577、1817、2745、11577；Y：496、4562、6912、9543。(1) X、Y 分别是什么元素？(2) X 与 Y 的最高价氧化物对应水化物能否反应？写出离子方程式。
\tcblower
\jie (1) X 的 $I_4\gg I_3$，最外层 3 个电子，为 Al；Y 的 $I_2\gg I_1$，最外层 1 个电子，$I_1$ 很小，为 Na。(2) 能，\ce{Al(OH)3} 是两性氢氧化物：\ce{Al(OH)3 + OH- = [Al(OH)4]-}。
\end{liti}

\subsection{电负性}
\begin{gainian}[电负性]
衡量元素原子在化合物中吸引键合电子能力的标度（F 最大，为 4.0）。同周期从左到右增大，同主族从上到下减小。
\end{gainian}
\textbf{应用：}①判断金属性、非金属性（金属一般 $<1.8$，非金属一般 $>1.8$）；②判断化合物中元素的正负化合价（电负性大的显负价，如 \ce{NaH} 中 H 为 $-1$ 价，\ce{OF2} 中 O 为 $+2$ 价）；③判断化学键类型（电负性差值大于 1.7 一般为离子键，小于 1.7 一般为共价键；差值越大，键的离子性越强）；④对角线规则（Li 与 Mg、Be 与 Al、B 与 Si 性质相似）。

\begin{liti}[电负性与化合价]
已知电负性：H 2.1，C 2.5，N 3.0，F 4.0，Si 1.8。判断下列化合物中加点元素的化合价：\ce{SiH4} 中的 H；\ce{NF3} 中的 N；\ce{CH2Cl2} 中的 C（Cl 为 3.0）。
\tcblower
\jie 电负性大的元素吸引电子能力强，显负价。\ce{SiH4}：H（2.1）$>$Si（1.8），H 为 $-1$ 价（\ce{SiH4} 中 Si 为 $+4$，有强还原性，遇水放出 \ce{H2}）。\ce{NF3}：F 为 $-1$，N 为 $+3$。\ce{CH2Cl2}：C 与 H 成键时 C 吸引电子（H 为 $+1$），与 Cl 成键时 Cl 吸引电子（Cl 为 $-1$），C 为 $0$ 价。
\end{liti}

\begin{yicuo}
第一电离能与电负性的大小顺序\textbf{不完全一致}：电负性 O$>$N，但第一电离能 N$>$O（N 的 $2p^3$ 半满）。比较金属性用电离能，比较非金属性用电负性。
\end{yicuo}

\section{元素推断的方法}

\begin{fangfa}[“位—构—性”三角推断]
\textbf{①由结构式推断}：根据原子形成的共价键数确定族（H、F、Cl 形成 1 个键；O、S 形成 2 个；N、P 形成 3 个，P 还可形成 5 个；C、Si 形成 4 个；B 形成 3 个，带负电时形成 4 个）。\textbf{②由电子排布推断}：如“3p 轨道半充满”为 P，“未成对电子数最多的第四周期元素”为 Cr，“$s$ 电子数与 $p$ 电子数相等”为 O 或 Mg。\textbf{③由电离能图推断}：找 ⅡA、ⅤA 的“凸起”和 0 族的“峰”、ⅠA 的“谷”。\textbf{④由性质推断}：地壳中含量最多的元素为 O，形成化合物种类最多的为 C，单质为淡黄色固体的为 S 等。
\end{fangfa}

\begin{zhenti}{2025年广东卷第9题\quad 来源：仓库所附真题文档}
元素 a～i 为短周期元素，原子序数依次为 $x\sim x+8$，其第一电离能与原子序数的关系为：从 a 到 f 总体增大，但 c 高于 d；f 为最高峰；g 骤降为最低；g、h、i 依次为低—稍高—稍低。下列说法正确的是
\xxb{a 和 g 同主族}{金属性：g$>$h$>$i}{原子半径：e$>$d$>$c}{最简单氢化物沸点：b$>$c}
\textbf{解析：}f 为峰、g 骤降，f 为稀有气体、g 为碱金属；c 比后一个元素 d 高，说明 c 为 ⅤA 族。连续 9 种元素中 f 为 Ne，则 a～i 依次为 B、C、N、O、F、Ne、Na、Mg、Al。验证：h（Mg）高于 i（Al），符合 ⅡA 特例。A：B 与 Na 不同族；B：Na、Mg、Al 金属性依次减弱，正确；C：同周期半径 c（N）$>$d（O）$>$e（F）；D：\ce{NH3} 分子间有氢键，沸点 \ce{NH3}$>$\ce{CH4}（c$>$b）。答案：B。
\end{zhenti}

\begin{zhenti}{2024年广东卷第12题\quad 来源：仓库所附真题文档}
一种可为运动员补充能量的物质，其分子结构式中，R、W、Z、X、Y 为原子序数依次增大的短周期主族元素，Z 和 Y 同族；结构式中 R 均形成 1 个键，W 形成 4 个键，Z 形成 3 个键，X 形成 2 个键，Y 形成 5 个键（含 1 个双键 X=Y）。则
\xxb{沸点：\ce{ZR3}$<$\ce{YR3}}{最高价氧化物的水化物的酸性：Z$<$W}{第一电离能：Z$<$X$<$W}{\ce{ZX3-} 和 \ce{WX3^2-} 空间结构均为平面三角形}
\textbf{解析：}Y 形成 5 个键且与 Z 同族，Z 形成 3 个键，Y 为 P、Z 为 N；W 形成 4 个键且序数小于 N，为 C；R 形成 1 个键、序数最小，为 H；X 形成 2 个键，序数在 N 与 P 之间，为 O。A：\ce{NH3} 分子间有氢键，沸点高于 \ce{PH3}，错；B：酸性 \ce{HNO3}$>$\ce{H2CO3}，错；C：第一电离能 C$<$O$<$N，错；D：\ce{NO3-}、\ce{CO3^2-} 中心原子价层电子对数均为 3，无孤电子对，平面三角形，正确（判断方法见第7章）。答案：D。
\end{zhenti}

% ---------------- 章末 ----------------
\zhangmo{本章小结}
\subsection*{一、知识总结}
\begin{center}
\small
\begin{tabularx}{\linewidth}{p{2.4cm}X}
\toprule
排布规则 & 构造原理（4s 先于 3d）、能量最低、泡利、洪特（含半满全满特例 Cr、Cu） \\
表示方法 & 电子排布式、简化式、价层排布式、轨道表示式；离子先失最外层电子 \\
周期表 & 周期$=$能层数；主族序数$=$价电子数；s、p、d、ds、f 五区 \\
周期律 & 半径：同周期减小、同主族增大；$I_1$：总体增大，ⅡA、ⅤA 特例；电负性：同周期增大、同主族减小 \\
\bottomrule
\end{tabularx}
\end{center}
\subsection*{二、方法提升}
\begin{enumerate}
\item 写排布式先按构造原理填，再按能层顺序整理；遇到 Cr、Cu 想到半满、全满。
\item 逐级电离能找“突跃”定最外层电子数；第一电离能比较注意 ⅡA、ⅤA。
\item 元素推断先抓最有特征的信息（成键数、半满、电离能峰谷），再用其他条件验证。
\end{enumerate}
\subsection*{三、易错点归纳}
\begin{enumerate}
\item \ce{Fe^2+} 写成 [Ar]$3d^44s^2$（错，应先失 4s）。
\item 认为第一电离能 O$>$N、Mg$<$Al（均错）。
\item 认为电负性大的元素第一电离能一定大。
\item 价层电子排布式与最外层电子排布式混淆（过渡元素价层包括 $(n-1)d$）。
\end{enumerate}

\zhangmo{分层练习}
\zu{A组\quad 基础巩固}
\begin{xiti}
\item 下列基态原子或离子的电子排布式（或轨道表示式）正确的是\xz
\xxa{Cr：[Ar]$3d^44s^2$}{\ce{Fe^2+}：[Ar]$3d^44s^2$}{Cu：[Ar]$3d^{10}4s^1$}{N 的 2p 轨道：\fbox{$\uparrow\downarrow$}\fbox{$\uparrow$}\fbox{\ }}

\item 下列各组元素第一电离能大小顺序正确的是\xz
\xx{Mg$<$Al$<$Si}{N$<$O$<$F}{Na$<$Mg$<$Al}{Al$<$Mg$<$P}

\item 下列说法正确的是\xz
\xxa{同一原子中，2p、3p、4p 能级的轨道数依次增多}{电子由 3s 能级跃迁到 3p 能级时释放能量}{s 区元素全部是金属元素}{焰色试验是电子跃迁时放出光能，属于物理变化}

\item 某元素的逐级电离能（kJ·mol$^{-1}$）依次为 738、1451、7733、10540、13630……，该元素最可能是\tk，其基态原子的价层电子排布式为\tk。

\item 下列说法正确的是\xz
\xxa{3d 能级最多容纳 6 个电子}{基态原子中 4s 能级的能量高于 3d 能级，所以先填 3d}{p 能级的 3 个原子轨道能量相同、相互垂直}{同一原子中，1s、2s、3s 电子的能量相同}

\item 写出下列粒子的价层电子排布式：\ce{_{24}Cr}\tk；\ce{_{29}Cu+}\tk；\ce{_{26}Fe^3+}\tk；\ce{_{35}Br}\tk。
\end{xiti}

\zu{B组\quad 高考典型}
\begin{xiti}
\item \laiyuan{2026年广东卷第11题，来源：仓库所附真题文档（网络收集版）}\ 化合物 \ce{L2MQR4} 与 \ce{TXR4} 可用作电解质饮料的配料。短周期元素 M、R、L、T、Q、X 原子序数依次增加，R、X 同族，Q 与 X 相邻，Q 基态原子的 3p 轨道为半充满，则下列说法正确的是\xz
\xxb{原子半径：R$>$X$>$M}{金属性：T$>$L$>$M}{第一电离能：Q$>$T$>$L}{最高价氧化物对应水化物的酸性：Q$>$X}

\item \laiyuan{原创}\ 前四周期元素 A、B、C、D、E 原子序数依次增大。A 的基态原子 2p 能级有 2 个未成对电子且 s 电子数与 p 电子数相等；B 是短周期中原子半径最大的主族元素；C 的单质是淡黄色固体；D 与 C 同周期，D 的电负性在同周期主族元素中最大；E 的基态原子 3d 能级有 6 个电子。(1) 写出 A～E 的元素符号。(2) 写出 E 的基态原子电子排布式和 \ce{E^3+} 的价层电子排布式，说明 \ce{E^3+} 较稳定的原因。(3) 比较 A、C、D 的第一电离能和电负性。(4) B 与 A 形成的两种化合物中，含非极性键的是\tk（写化学式）。

\item \laiyuan{原创}\ 短周期主族元素 X、Y、Z、W 原子序数依次增大。X 的基态原子 2p 能级上有 3 个未成对电子；Y 是地壳中含量最多的元素；Z 的第一电离能在同周期主族元素中最小；W 与 Y 同主族。下列说法正确的是\xz
\xxa{第一电离能：Y$>$X$>$W}{电负性：Y$>$W$>$X}{简单离子半径：W$>$Y$>$Z}{Z 与 Y 形成的化合物中只含离子键}

\item \laiyuan{原创}\ 解释下列事实：(1) 第一电离能 Mg$>$Al；(2) 第一电离能 P$>$S；(3) 离子半径 $r(\ce{S^2-})>r(\ce{Cl-})>r(\ce{K+})>r(\ce{Ca^2+})$；(4) Na 的第二电离能远大于其第一电离能。
\end{xiti}

\zu{C组\quad 冲刺高分}
\begin{xiti}
\item \laiyuan{原创}\ (1) 基态 \ce{Cr} 原子的价层轨道表示式为\tk，其未成对电子数为\tk。(2) 比较 Mn 和 Fe 的第三电离能 $I_3$ 的大小并解释原因。(3) 第二周期元素中，第一电离能介于 B 和 N 之间的元素有\tk{}种。(4) 已知 Cu 的 $I_2=1958$ kJ·mol$^{-1}$，Zn 的 $I_2=1733$ kJ·mol$^{-1}$，解释为何 $I_2(\ce{Cu})>I_2(\ce{Zn})$。

\item \laiyuan{原创}\ 锂离子电池正极材料常含 Fe、Co、Ni、Mn 等过渡元素。(1) 写出基态 \ce{Co^3+} 和 \ce{Ni^2+} 的价层电子排布式，并指出各自的未成对电子数。(2) \ce{Mn^2+} 很难被氧化为 \ce{Mn^3+}，而 \ce{Fe^2+} 较易被氧化为 \ce{Fe^3+}，从电子排布角度解释。(3) Co 位于周期表第\tk{}周期第\tk{}族，属于\tk{}区。
\end{xiti}

\zhangmo{练习解析}
\begin{jiexilist}
\item \ans{C}A：Cr 为 [Ar]$3d^54s^1$（半满稳定）；B：\ce{Fe^2+} 为 [Ar]$3d^6$（先失 4s）；C：Cu 为 $3d^{10}4s^1$（全满稳定），正确；D：违反洪特规则，2p 的 3 个电子应分占 3 个轨道且自旋平行。
\item \ans{D}A：Mg（$3s^2$ 全满）$>$Al，应为 Al$<$Mg$<$Si；B：N（$2p^3$ 半满）$>$O，应为 O$<$N$<$F；C：应为 Na$<$Al$<$Mg；D：Al$<$Mg$<$P，正确。
\item \ans{D}A：p 能级都有 3 个轨道；B：3s$\to$3p 是向高能级跃迁，吸收能量；C：H 在 s 区，是非金属；D 正确。
\item Mg；$3s^2$。$I_3\gg I_2$，最外层有 2 个电子，为 ⅡA 族；$I_1=738$ 较小，为第三周期的 Mg（Be 的 $I_1$ 约 900）。
\item \ans{C}A：d 能级有 5 个轨道，最多 10 个电子；B：构造原理中 4s 能量低于 3d，先填 4s；C 正确；D：能层越高能量越高，$E(1s)<E(2s)<E(3s)$。
\item $3d^54s^1$；$3d^{10}$；$3d^5$；$4s^24p^5$。Cr 半满稳定；\ce{Cu+} 由 $3d^{10}4s^1$ 失去 4s 电子；\ce{Fe^3+} 先失 2 个 4s 电子再失 1 个 3d 电子；主族元素的价层电子即最外层电子。
\item \ans{C}Q 的 3p 半充满，为 P；X 与 Q 相邻且在其后，为 S；R 与 X 同族且序数最小，为 O；化合物 \ce{L2MQR4}（如 \ce{K2HPO4} 类）和 \ce{TXR4}（如 \ce{MgSO4}）中，M 序数最小可形成 \ce{L2MPO4}，M 为 H（\ce{Na2HPO4}），L 为 Na，T 为 Mg。验证：原子序数 H$<$O$<$Na$<$Mg$<$P$<$S。A：半径 S$>$O$>$H，即 X$>$R$>$M，错；B：金属性 Na$>$Mg，即 L$>$T，错；C：第一电离能 P$>$Mg$>$Na，即 Q$>$T$>$L，正确；D：酸性 \ce{H2SO4}$>$\ce{H3PO4}，即 X$>$Q，错。答案：C。
\item (1) A 为 O（$1s^22s^22p^4$，s、p 电子均为 4，2p 有 2 个未成对电子）；B 为 Na；C 为 S；D 为 Cl；E 为 Fe（$3d^64s^2$）。(2) Fe：$1s^22s^22p^63s^23p^63d^64s^2$；\ce{Fe^3+}：$3d^5$，3d 能级半充满，能量较低，较稳定。(3) 第一电离能：Cl$>$O$>$S（同周期 S$<$Cl，同主族 O$>$S；O 与 Cl 比较：O 为 1314、Cl 为 1251 kJ·mol$^{-1}$，O$>$Cl）——高考一般只考查同周期或同主族的比较；电负性：O$>$Cl$>$S。(4) \ce{Na2O2}（含 O—O 非极性键）。
\item \ans{C}X 为 N（$2p^3$），Y 为 O，Z 为 Na（同周期 $I_1$ 最小），W 为 S。A：第一电离能 N$>$O$>$S，即 X$>$Y$>$W，错；B：电负性 O$>$N$>$S，即 Y$>$X$>$W，错；C：\ce{S^2-} 有 3 个电子层，半径最大；\ce{O^2-} 与 \ce{Na+} 电子层结构相同，核电荷数小的半径大，\ce{O^2-}$>$\ce{Na+}，正确；D：\ce{Na2O2} 中含 O—O 共价键，错。
\item (1) Mg 的价层电子为 $3s^2$（全满，较稳定），Al 为 $3s^23p^1$，3p 电子能量较高且易失去，所以 Al 的 $I_1$ 较小。(2) P 的 3p 为 $3p^3$ 半满，较稳定；S 为 $3p^4$，失去一个电子后变为半满，较容易。(3) 四种离子电子层结构相同（均为 Ar 型），核电荷数依次增大，核对电子的吸引依次增强，半径依次减小。(4) Na 失去 1 个电子后变为 \ce{Na+}（$2s^22p^6$，稀有气体稳定结构），再失去的电子位于内层，离核近、受吸引强，所以 $I_2\gg I_1$。
\item (1) 3d：\fbox{$\uparrow$}\fbox{$\uparrow$}\fbox{$\uparrow$}\fbox{$\uparrow$}\fbox{$\uparrow$}\ 4s：\fbox{$\uparrow$}；6。(2) $I_3(\ce{Mn})>I_3(\ce{Fe})$。\ce{Mn^2+} 为 $3d^5$ 半满，较稳定，再失去一个电子较难；\ce{Fe^2+} 为 $3d^6$，失去一个电子后变为较稳定的 $3d^5$，较容易。(3) 3 种：Be、C、O（第一电离能 B$<$Be$<$C$<$O$<$N）。(4) Cu 失去 1 个电子后 \ce{Cu+} 为 $3d^{10}$ 全满稳定结构，再失电子较难；Zn 失去 1 个电子后 \ce{Zn+} 为 $3d^{10}4s^1$，再失去 4s 电子较容易。
\item (1) \ce{Co^3+}：$3d^6$，4 个未成对电子；\ce{Ni^2+}：$3d^8$，2 个未成对电子。(2) \ce{Mn^2+} 为 $3d^5$ 半满稳定结构，失电子后破坏半满，较难；\ce{Fe^2+} 为 $3d^6$，失去 1 个电子后变为 $3d^5$ 半满稳定结构，较容易。(3) 四；Ⅷ；d。
\end{jiexilist}
